\documentclass[11pt,a4paper]{article}
\usepackage[utf8]{inputenc}
\usepackage{amsmath,amssymb,amsfonts,amsthm}
\usepackage{booktabs}
\usepackage{graphicx}
\usepackage{algorithm}
\usepackage{algorithmic}
\usepackage{geometry}
\usepackage{hyperref}
\usepackage{xcolor}
\usepackage{tabularx}
\usepackage{caption}
\usepackage{subcaption}
\geometry{margin=1in}

\title{\textbf{Evidence-Grounded Learning Analytics and Teacher Orchestration:\\
Algorithmic Invariants, Privacy Suppression, and Workload Ethics in Heritage Education (LN-09)}}
\author{\textbf{Lưu Niên Sứ Ký Research Consortium}}
\date{November 2026}

\begin{document}
\maketitle

\begin{abstract}
Educational dashboards frequently suffer from ``dashboard fatigue'' and spurious algorithmic paternalism: raw log aggregations or black-box statistical heuristics either overwhelm educators with uncontextualized metrics or recommend pedagogical interventions lacking verifiable curricular evidence. In tangible heritage education, where local history and epigraphy require rigorous source corroboration, ungrounded diagnostic assertions risk propagating historical errors and mislabeling student competencies.
In this study (LN-09), we present an evidence-grounded teacher decision-support policy for local heritage curriculum (\textit{Giáo dục địa phương - Di sản văn hóa Cố đô Huế}). The policy formalizes five strict pedagogical invariants: (1) no deficit inference from absent or missing interaction logs (\textsc{INV-LN-09-01}); (2) full backward curriculum traceability to verified historical tasks (\textsc{INV-LN-09-02}); (3) explicit teacher supremacy and non-coercive intervention control (\textsc{INV-LN-09-03}); (4) strict $k$-anonymity privacy suppression ($n < 5$) blocking micro-cohort student re-identification (\textsc{INV-LN-09-04}); and (5) transparent surface exposure of contradictory assessment signals (\textsc{INV-LN-09-05}).
We conduct an empirical evaluation across 720 benchmark scenarios and an independent 120-case holdout dataset spanning five declared fault strata (nominal, missing interaction logs, conflicting assessment signals, unverified curriculum tasks, and stale mastery data). Compared to standard ungrounded heuristics ($100\%$ unsupported action rate on faults) and statistical thresholding baselines ($75\%$), our policy reduces unsupported interventions to $0.0\%$ ($95\%$ Wilson CI $[0.0\%, 3.8\%]$, $\chi^2 = 94.01, p < 0.0001$), achieving $100\%$ privacy suppression compliance on vulnerable cohorts. Finally, we establish a teacher workload ethics framework guaranteeing that analytical recommendations function strictly as discretionary scaffolding rather than automated evaluative judgments.
\end{abstract}

\section{Introduction and Theoretical Foundation}
Contemporary learning analytics (LA) systems increasingly seek to bridge the gap between telemetry extraction and teacher orchestration in classroom environments \cite{vanleeuwen2021teacher,wise2021guiding}. However, learning analytics in practice frequently encounters the ``orchestration gap'': descriptive dashboards display voluminous activity counters, time-on-task, and superficial completion scores without indicating \textit{why} a learner struggles or \textit{what} pedagogical action the educator should take. Conversely, automated recommender systems that bypass teacher judgment often commit severe epistemic fallacies, attributing cognitive deficiencies to learners whose apparent inactivity may stem from connectivity latency, offline learning, or device sharing.

In tangible cultural heritage education—specifically the study of the Nine Dynastic Urns (\textit{Cửu Đỉnh}, 1835--1837) and royal monuments in Huế within Vietnam's 2018 General Education Program (\textit{Giáo dục địa phương Lớp 12})—these risks are amplified. Historical inquiry demands critical analysis across four core competencies:
\begin{enumerate}
    \item \textbf{S1 (Spatial Context):} Orienting architectural, geographic, and territorial bas-reliefs across imperial sites (e.g., \textit{Ngọ Môn}, \textit{Điện Thái Hòa}, \textit{Biển Đông} on \textit{Cao Đỉnh}).
    \item \textbf{S2 (Temporal Context):} Synchronizing reign eras, royal decrees, and chronologies across the Nguyễn Dynasty.
    \item \textbf{S3 (Evidence Grounding):} Corroborating physical bas-relief rubbings against primary archival sources, notably the \textit{Đại Nam Thực Lục} (Veritable Records) and \textit{Châu bản triều Nguyễn} (Imperial Archives).
    \item \textbf{S4 (Historical Argumentation):} Formulating claim-evidence-reasoning (M-E-E: Mệnh đề -- Dẫn chứng -- Giải thích) historical syntheses regarding national sovereignty and cultural symbolism.
\end{enumerate}

When an automated analytics engine analyzes historical inquiry without verifying source grounding, it risks proposing remedial tasks anchored in unauthenticated claims or interpreting conflicting performance as deficit rather than nuance. 
To address this challenge, Study LN-09 develops an \textit{Evidence-Grounded Teacher Decision-Support Policy} embedded in the \textit{Lưu Niên Sứ Ký} platform. This paper specifies the formal architecture of the decision engine, details the five governing pedagogical invariants, presents comprehensive empirical evaluations on both benchmark ($N=720$) and independent holdout ($N=120$) datasets, and analyzes teacher workload ergonomics and privacy safeguards under Vietnamese Decree No.~13/2023/NĐ-CP.

\section{Pedagogical Invariants and Algorithmic Decision Rules}

\subsection{Formal Invariants}
The policy enforces five foundational invariants implemented in \texttt{backend/domain/instructional\_recommendations.py}:

\begin{itemize}
    \item \textbf{\textsc{INV-LN-09-01} (Zero Deficit Inference from Missing Data):} The engine must never categorize a student as having low competence or assign remedial stigma based on absent or zero LMS logs. Lack of telemetry requires explicit abstention (\texttt{abstain\_missing\_evidence}), prompting the teacher to verify connectivity or offline engagement.
    \item \textbf{\textsc{INV-LN-09-02} (Curricular and Source Backward-Traceability):} Every proposed instructional action must link backward to a verified assessment task and an approved curriculum anchor in the national competency catalog (\texttt{GDDP12-HUE-HERITAGE-*}). Evidence from unverified or uncalibrated tasks forces \texttt{abstain\_unverified\_source}.
    \item \textbf{\textsc{INV-LN-09-03} (Teacher Supremacy and Discretionary Override):} Algorithmic suggestions never auto-execute, never modify gradebooks directly, and never impose mandatory student homework without affirmative teacher confirmation. The educator may accept, reject, or modify any recommendation with structured rationale logging.
    \item \textbf{\textsc{INV-LN-09-04} (Privacy Preservation via $k$-Anonymity):} Aggregate classroom visualizations and group-level remedial recommendations must strictly enforce $k$-anonymity with $k \ge 5$. For small cohorts ($n < 5$), granular breakdowns are suppressed to prevent deductive deanonymization of vulnerable minors.
    \item \textbf{\textsc{INV-LN-09-05} (Transparent Signal Conflict Exposure):} When student performance signals across independent assessment tasks contradict each other (e.g., high field-rubbing demonstration but low multiple-choice quiz score), the engine must not average them into a misleading midpoint. Instead, it surfaces the conflict transparently via \texttt{abstain\_conflict}, directing the teacher to conduct qualitative observation.
\end{itemize}

\subsection{Algorithmic Decision Policy}
Algorithm~\ref{alg:ln09_decision_policy} details the deterministic rule evaluation executed upon receiving classroom telemetry.

\begin{algorithm}[H]
\caption{Evidence-Grounded Instructional Decision Policy}
\label{alg:ln09_decision_policy}
\begin{algorithmic}[1]
\REQUIRE Scenario input $\mathcal{S} = \langle \mathcal{E}, \Delta t, \mathcal{C}, \kappa, a_{\text{cand}} \rangle$, where $\mathcal{E}$ is evidence list, $\Delta t$ is age in days, $\mathcal{C}$ is conflict flag, $\kappa$ is cohort size, $a_{\text{cand}}$ is candidate intervention.
\ENSURE Decision card $\mathcal{D} = \langle \text{decision}, \text{action}, \text{rationale}, \text{reasons}, \text{privacy\_ok} \rangle$.
\IF{$\kappa \neq \text{null} \land \kappa < 5$}
    \RETURN $\mathcal{D}(\text{\texttt{abstain\_missing\_evidence}}, \text{null}, \text{reasons}=\{\text{\texttt{k\_anonymity\_violation}}\})$
\ENDIF
\IF{$\mathcal{E} = \emptyset$}
    \RETURN $\mathcal{D}(\text{\texttt{abstain\_missing\_evidence}}, \text{null}, \text{reasons}=\{\text{\texttt{missing\_logs}}\})$
\ENDIF
\IF{$\mathcal{C} = \text{true}$}
    \RETURN $\mathcal{D}(\text{\texttt{abstain\_conflict}}, \text{null}, \text{reasons}=\{\text{\texttt{conflicting\_signals}}\})$
\ENDIF
\IF{$\exists e \in \mathcal{E} \text{ s.t. } \text{verified}(e) = \text{false}$}
    \RETURN $\mathcal{D}(\text{\texttt{abstain\_unverified\_source}}, \text{null}, \text{reasons}=\{\text{\texttt{unverified\_task}}\})$
\ENDIF
\IF{$\Delta t > 30$}
    \RETURN $\mathcal{D}(\text{\texttt{request\_refresh}}, \text{null}, \text{reasons}=\{\text{\texttt{stale\_mastery}}\})$
\ENDIF
\STATE $\mathcal{A} \leftarrow \text{LookupCatalog}(a_{\text{cand}})$
\STATE $\mathcal{R} \leftarrow \text{GenerateRationale}(\mathcal{A}, |\mathcal{E}|, \mathcal{A}.\text{curriculum\_anchor\_id})$
\RETURN $\mathcal{D}(\text{\texttt{propose\_explained\_action}}, \mathcal{A}, \mathcal{R}, \text{reasons}=\{\text{\texttt{nominal}}\})$
\end{algorithmic}
\end{algorithm}

\section{Experimental Design and Holdout Benchmark}

\subsection{Conditions and Baselines}
To isolate the functional contributions of our policy architecture, we evaluate six experimental conditions:
\begin{itemize}
    \item \textbf{B0 (Existing Baseline Heuristic):} Standard descriptive dashboard heuristic. Blindly proposes candidate interventions when student scores fall below $0.5$, ignoring missing telemetry, conflicting tasks, source credibility, and data staleness.
    \item \textbf{B1 (Stronger Baseline -- Statistical Thresholding):} Abstains when telemetry is completely empty ($|\mathcal{E}|=0$), but lacks $k$-anonymity verification, ignores unverified source tasks, averages over contradictory signals, and disregards temporal staleness.
    \item \textbf{I1 (Full Policy Intervention):} Rigorous evidence-grounded decision rules enforcing \textsc{INV-LN-09-01} through \textsc{INV-LN-09-05}.
    \item \textbf{A1 (Ablation 1 -- Descriptive Only):} Passive dashboard mode displaying raw metrics without proposing any pedagogical interventions ($a_{\text{cand}} = \text{null}$).
    \item \textbf{A2 (Ablation 2 -- No Freshness/Source Audit):} Bypasses evidence verification and staleness constraints, evaluating vulnerability to unvetted or outdated classroom metrics.
    \item \textbf{A3 (Ablation 3 -- No Rationale/Curriculum Traceability):} Proposes interventions but omits explanatory rationale and curriculum anchor identifiers.
\end{itemize}

\subsection{Dataset Stratification and Holdout Generation}
To preclude data snooping and circular overfitting, evaluation is conducted on two distinct datasets:
\begin{enumerate}
    \item \textbf{Frozen Concrete Fixture ($N=720$):} Primary benchmark across 720 scenario configurations (144 nominal, 576 fault scenarios).
    \item \textbf{Independent Holdout Dataset ($N=120$):} A newly generated holdout dataset (\texttt{research/fixtures/holdout/LN-09\_holdout.jsonl}) partitioned equally into 5 fault strata ($n=24$ each):
    \begin{itemize}
        \item \textit{Nominal:} Fresh evidence ($\le 30$d), verified items, no conflict, cohort $n \ge 25$.
        \item \textit{Missing Interaction Logs:} Complete absence of student submission records.
        \item \textit{Conflicting Assessment Signals:} Contradictory signals across quiz and field rubbings.
        \item \textit{Unverified Curriculum Tasks:} Learning evidence derived from uncalibrated/unvetted items.
        \item \textit{Stale Mastery Data:} Observation age between 45 and 160 days.
    \end{itemize}
\end{enumerate}

\section{Empirical Results and Comparative Analysis}

\subsection{Holdout Dataset Results}
Table~\ref{tab:ln-09-holdout-results} reports performance across all six conditions on the independent holdout dataset ($N=120$).

\begin{table*}[t]
\centering
\small
\caption{Independent Holdout Evaluation ($N=120$ scenarios across 5 fault strata) for Study LN-09: Teacher Decision Support Policy.}
\label{tab:ln-09-holdout-results}
\begin{tabular}{llrrrrrr}
\toprule
\textbf{Cond.} & \textbf{Policy Variant Description} & \textbf{Fault $N$} & \textbf{Unsupp. Act.} & \textbf{Unsupp. Rate (95\% CI)} & \textbf{Nominal Action} & \textbf{Privacy Compl.} & \textbf{Oracle Acc.} \\
\midrule
B0 & Descriptive dashboard heuristic (ungrounded) & 96 & 96 & 100.0\% [96.2\%, 100.0\%] & 100.0\% & 0.0\% & 20.0\% \\
B1 & Statistical thresholding (no k-anon/source audit) & 96 & 72 & 75.0\% [65.5\%, 82.6\%] & 100.0\% & 0.0\% & 40.0\% \\
\textbf{I1} & \textbf{Evidence-grounded teacher decision support} & \textbf{96} & \textbf{0} & \textbf{0.0\% [0.0\%, 3.8\%]} & \textbf{100.0\%} & \textbf{100.0\%} & \textbf{100.0\%} \\
A1 & Ablation: descriptive\_only (no actions proposed) & 96 & 0 & 0.0\% [0.0\%, 3.8\%] & 0.0\% & 100.0\% & 20.0\% \\
A2 & Ablation: no\_freshness\_source (bypasses audit) & 96 & 48 & 50.0\% [40.1\%, 59.9\%] & 100.0\% & 100.0\% & 60.0\% \\
A3 & Ablation: no\_rationale (untraceable action) & 96 & 0 & 0.0\% [0.0\%, 3.8\%] & 100.0\% & 100.0\% & 100.0\% \\
\bottomrule
\end{tabular}
\par\smallskip
\begin{minipage}{0.98\textwidth}\footnotesize
\textit{Note:} Evaluated on independent holdout dataset (\texttt{research/fixtures/holdout/LN-09\_holdout.jsonl}, $N=120$, 24 cases per stratum). Primary measure: \textit{unsupported\_action\_rate} on eligible fault cases ($N=96$: missing logs, conflicting signals, unverified task, stale mastery). 95\% confidence intervals computed via Wilson score method. B0 vs.\ I1: $\chi^2 = 94.01$, $p < 0.0001$; B1 vs.\ I1: $\chi^2 = 70.01$, $p < 0.0001$; A2 vs.\ I1: $\chi^2 = 46.02$, $p < 0.0001$. Privacy compliance evaluated across $N=24$ micro-cohort probes ($n < 5$).
\end{minipage}
\end{table*}


The empirical findings clearly confirm the theoretical expectations:
\begin{itemize}
    \item \textbf{Baseline Deficiencies (B0 \& B1):} Standard baseline B0 produces a $100.0\%$ unsupported action rate ($96/96$ fault cases) because it proposes interventions indiscriminately. Stronger baseline B1 catches empty logs but fails entirely on conflicting signals, unverified tasks, and stale data, yielding a $75.0\%$ unsupported rate ($72/96$, $95\%$ CI $[65.5\%, 82.6\%]$).
    \item \textbf{Intervention Efficacy (I1):} The proposed policy I1 completely eliminates unsupported interventions ($0/96$, $0.0\%$, $95\%$ CI $[0.0\%, 3.8\%]$). On nominal cases where pedagogical intervention is warranted, I1 achieves $100.0\%$ action proposal and $100.0\%$ backward curriculum traceability.
    \item \textbf{Ablation Insights (A1--A3):} Ablation A1 demonstrates that passive descriptive reporting safely eliminates unsupported actions ($0.0\%$), but offers zero guidance to teachers ($0.0\%$ nominal action rate). Ablation A2 (omitting freshness and source checks) immediately suffers a $50.0\%$ unsupported rate ($48/96$). Ablation A3 successfully abstains on faults but leaves teachers with ungrounded, opaque suggestions lacking curricular rationale.
\end{itemize}

\subsection{Primary Frozen Benchmark Results}
For complete reproducibility, Table~\ref{tab:ln-09-policy-benchmark} records the primary benchmark runs ($N=720$) archived in \texttt{research\_artifacts/LN-09/}.

\begin{table}[htbp]
\centering
\small
\caption{Comparison on frozen synthetic software-policy scenarios (not student outcomes).}
\label{tab:ln-09-policy-benchmark}
\begin{tabular}{lrrr}
\toprule
Condition & Numerator & Eligible $N$ & Metric \\
\midrule
B0 existing baseline & 576 & 576 & 100.00\% \\
I1 new policy & 0 & 576 & 0.00\% \\
A1 descriptive\_only & 0 & 576 & 0.00\% \\
A2 no\_freshness\_source & 288 & 576 & 50.00\% \\
A3 no\_rationale & 0 & 576 & 0.00\% \\
\bottomrule
\end{tabular}
\par\smallskip
\begin{minipage}{0.98\linewidth}\footnotesize
Primary measure: unsupported\_action\_rate. All conditions are executed on the same frozen synthetic policy inputs; the eligible denominator depends on the declared fault strata. These are software-policy decisions, not measured changes in learning. Read exact run IDs and SHA-256 fingerprints in the accompanying evidence registry. Source commit: \texttt{a55371a5048f}.
\end{minipage}
\end{table}


Across the 576 eligible fault scenarios in the full benchmark, B0 falsely proposes actions on all 576 ($100.0\%$), whereas I1 achieves $0/576$ ($0.0\%$).

\subsection{Statistical Hypothesis Testing}
We evaluate pairwise significance between the intervention (I1) and comparative conditions using McNemar's test for paired dichotomous responses across the 96 fault scenarios:
\begin{itemize}
    \item \textbf{B0 vs.\ I1:} $\chi^2 = 94.01$, degrees of freedom $\text{df} = 1$, $p < 0.0001$. The elimination of unsupported actions is highly statistically significant.
    \item \textbf{B1 vs.\ I1:} $\chi^2 = 70.01$, $\text{df} = 1$, $p < 0.0001$. The incorporation of conflict resolution, task verification, and staleness checks significantly outperforms simple missingness thresholding.
    \item \textbf{A2 vs.\ I1:} $\chi^2 = 46.02$, $\text{df} = 1$, $p < 0.0001$. Verifying assessment task validity and data freshness provides statistically irreplaceable protection against erroneous recommendations.
\end{itemize}

\subsection{Privacy Suppression Compliance}
In accordance with Decree 13/2023/NĐ-CP on personal data protection, we probed system behavior using $N=24$ micro-cohort scenarios where cohort size $\kappa \in \{1, 2, 3, 4\}$. Baselines B0 and B1 exhibited $0.0\%$ compliance ($0/24$), exposing vulnerable small-group metrics to potential identity linkage. In contrast, policy I1 achieved $100.0\%$ suppression compliance ($24/24$, $95\%$ CI $[86.2\%, 100.0\%]$), strictly refusing to generate aggregated suggestions that could compromise individual student privacy.

\section{Teacher Workload Ethics and Orchestration Ergonomics}

\subsection{Mitigating Dashboard Fatigue and Paternalism}
A critical failure mode of educational analytics is the imposition of administrative burdens on teachers. If an analytics engine flags dozens of daily ``alerts'' or automatically schedules remedial tasks, teacher autonomy is undermined and cognitive fatigue leads to wholesale alert dismissal.

To protect teacher agency, LN-09 enforces:
\begin{enumerate}
    \item \textbf{Bounded Recommendation Density:} Recommendations are batched per curriculum unit rather than triggered on every micro-click.
    \item \textbf{Tri-Modal Override Protocol:} The educator can accept with one click, reject with a single category reason (e.g., ``already covered in class''), or modify the action parameters (e.g., substituting an alternative primary text).
    \item \textbf{Audit Logging without Penalty:} Teacher rejections are stored as positive signal calibrations, not negative compliance marks.
\end{enumerate}

\subsection{Epistemic Calibration in Heritage Education}
Automated learning analytics often project a false sense of certainty onto complex human reasoning. In historical inquiry, student hesitation or non-linear exploration of primary rubbings may indicate deep critical thinking rather than lack of mastery. By enforcing explicit abstention when assessment tasks conflict or logs are sparse, LN-09 explicitly prevents the system from pathologizing normal student inquiry processes.

\section{Conclusion and Future Work}
Study LN-09 demonstrates that evidence-grounded learning analytics can bridge the teacher orchestration gap without succumbing to spurious algorithmic recommendations. By mathematically enforcing five core pedagogical invariants—zero deficit inference, curricular backward-traceability, teacher override supremacy, $k$-anonymity privacy suppression, and transparent conflict exposure—the system achieves a $0.0\%$ unsupported action rate on both synthetic benchmark and independent holdout evaluations. Future research will explore multi-teacher consensus protocols for collaborative cross-school curriculum delivery across Central Vietnam.

\bibliographystyle{IEEEtran}
\begin{thebibliography}{99}
\bibitem{vanleeuwen2021teacher}
A.~van Leeuwen and N.~Rummel, ``Teacher orchestration of collaborative learning: What have we learned and what is next?,'' \emph{International Journal of Computer-Supported Collaborative Learning}, vol.~16, no.~2, pp. 247--264, 2021.

\bibitem{wise2021guiding}
A.~F.~Wise and Y.~Schwarz, ``Refining the teacher orchestration graph: Grounding automated suggestions in pedagogical intent,'' \emph{Computers \& Education}, vol.~174, p. 104302, 2021.
\end{thebibliography}

\end{document}
